% !TeX root = ../main.tex
\section{Задание и исходные данные}

Целью данной работы является исследование алгоритмов пространственно-временной обработки сигналов на примере устройства цифровой фильтрации радиосигналов, включающего аналого-цифровой преобразователь (АЦП), цифровой гетеродин (ЦГ) и каскад децимирующих фильтров.

В соответствии с заданием, для варианта №8 ($N_{var} = 8$) параметры разрабатываемого устройства приведены в таблицу \ref{tab:params}
\begin{table}[H]
    \centering
    \caption{Исходные параметры для варианта 8}
    \label{tab:params}
    \vspace{2mm}
    \begin{tabular}{|l|c|c|}
        \hline
        \textbf{Наименование параметра} & \textbf{Формула} & \textbf{Значение} \\ \hline
        Промежуточная частота входного сигнала, $f_c$ & $130 + N_{var} \cdot 10$ & $210$ МГц \\ \hline
        Количество бит АЦП, $n_a$ & $5 + N_{var}$ & $13$ бит \\ \hline
        Частота дискретизации АЦП, $f_s$ & $100 + N_{var} \cdot 10$ & $180$ МГц \\ \hline
        Количество бит ЦГ, $n_g$ & $5 + N_{var}$ & $13$ бит \\ \hline
        Центральная частота ЦГ, $f_G$ & $130 + N_{var} \cdot 10$ & $210$ МГц \\ \hline
        База ЛЧМ-сигнала, $B$ & - & $512$ \\ \hline
        Коэффициент децимации (из скрипта), $R$ & - & $10$ \\ \hline
    \end{tabular}
\end{table}

